% ===== supplementary/hoare_ltl_formalism.tex =====
% Archived from sec/07_methodology.tex (Hoare-triple and LTL subsections).
% Removed from main paper per rewrite plan item 7 (Hoare) and 8 (LTL).

\section{Hoare-Triple Obligation per Step}
\label{sec:hoare-archived}

Each step of Algorithm~\ref{alg:enum} carries a Hoare-triple
obligation in the sense of
Hoare~\cite{hoare1969axiomatic} and
Dijkstra~\cite{dijkstra1975guarded}:
\[
\{\,P_t\,\}\quad w_t\quad\{\,P_{t+1}\,\}
\]
where $P_t$ is the precondition holding on the state space
before the $t$-th emitted operation $w_t$, and $P_{t+1}$ is the
postcondition that the verifier checks via its
$\delta_\mathcal{V}$ transition.  When the precondition fails,
the verifier rejects; when the postcondition fails, the verifier
also rejects; the trace is admitted iff every Hoare triple
along the emitted path holds.

This is the formal sense in which the architecture is
\emph{validation-by-construction}: the verifier's DFA is the
mechanical realisation of a finite Hoare-style proof obligation
for the Goal-Driven generator's emitted trace.

\section{Temporal-Logic Statement of Correctness}

The correctness property of the binding is captured in
LTL~\cite{pnueli1977temporal} by the safety-and-liveness pair
\begin{align*}
\textsf{Safety:}\quad &
G\bigl(\,\mathsf{emit}\,\to\,\mathsf{verifier-tracking}\bigr) \\
\textsf{Liveness:}\quad &
G\bigl(\,\mathsf{plan}\,\to\,F\,(\mathsf{verified}\,\vee\,
\mathsf{halt-on-tape-bound})\bigr).
\end{align*}
Safety says that every emitted symbol advances the verifier's
state; liveness says that every plan eventually either succeeds
verification or hits the tape bound.  Both can be checked on the finite orchestrator control-flow graph
used in the implementation~\cite{clarke1981ctl,clarke1999modelchecking}.
